% Compile from the repository root:
% latexmk -pdf -interaction=nonstopmode -halt-on-error \
%   -outdir=/tmp/ito-covariance-note \
%   docs/ito_tracers_multidimensional_covariance_note_for_eric_moseley.tex
\documentclass[11pt]{article}

\usepackage[T1]{fontenc}
\usepackage{amsmath}
\usepackage{amssymb}
\usepackage{booktabs}
\usepackage{geometry}
\usepackage{xcolor}
\usepackage{hyperref}
\usepackage{microtype}

\geometry{margin=0.9in}
\hypersetup{
  colorlinks=true,
  urlcolor=blue!55!black,
  pdftitle={A finite-timestep covariance term in multidimensional Ito-2},
  pdfauthor={Drummond Fielding}
}
\setlength{\parindent}{0pt}
\setlength{\parskip}{0.55em}

\newcommand{\ddX}{\Delta\mathbf{X}}
\newcommand{\mean}[1]{\mathbb{E}\!\left[#1\right]}
\newcommand{\cov}{\operatorname{Cov}}
\newcommand{\diag}{\operatorname{diag}}

\title{\textbf{A finite-timestep covariance term\\
in multidimensional It\^o-2}}
\author{}
\date{}

\begin{document}
\maketitle
\vspace{-2.5em}

Eric,

I think there is one small but real finite-\(\Delta t\) issue in the
multidimensional argument. Nothing changes in 1D. It also disappears in the
strict \(\Delta t\rightarrow0\) limit. But for the actual one-kick-per-hydro-step
algorithm, independent coordinate kicks match the marginal variances, not the
full MC covariance.

\section*{Exact spot in the paper}

I mean Section 2.2 of
\href{https://arxiv.org/abs/2604.23041v2}{arXiv:2604.23041v2},
printed page 7.

The 1D result is fine. Equations (28)--(29) give
\begin{equation}
  u_i=\frac{h_i}{\Delta t}C_{-,i},
  \qquad
  \kappa_{ii}
  =
  \frac{h_i^2}{2\Delta t}
  \left(C_{+,i}-C_{-,i}^2\right).
\end{equation}
The \(-C_{-,i}^2\) term is the subtraction of the squared mean.

The issue starts at equations (45)--(46), where the paper calls
\begin{equation}
  \kappa_{ij}
  =
  \frac{\mean{\Delta X_i\Delta X_j}}{2\Delta t}
\end{equation}
the covariance matrix and then uses
\begin{equation}
  \mean{\Delta X_i\Delta X_j}
  =
  \mean{\Delta X_i^2}\delta_{ij}.
\end{equation}
The second equation is right for the \emph{raw} second moment.
For a finite-step covariance, the matching generalization of equation (29) is
\begin{equation}
  \boxed{
  \kappa_{ij}^{(\mathrm{step})}
  =
  \frac{
  h_i^2 C_{+,i}\delta_{ij}
  -
  h_i h_j C_{-,i}C_{-,j}
  }{2\Delta t}.
  }
\end{equation}

So the sentence immediately after equation (46), saying that the matrix is
diagonal and the dimensions are independent, is only true for the raw second
moment, or in the strict infinitesimal-rate limit.

Equations (47)--(48) can stay. Equation (49) should be written
\begin{equation}
  \boxed{
  \boldsymbol{\Sigma}\boldsymbol{\Sigma}^{\mathsf T}
  =
  2\boldsymbol{\kappa},
  }
\end{equation}
not as an element-by-element square root, unless
\(\boldsymbol{\kappa}\) is diagonal.

Equation (59) can also stay:
\begin{equation}
  \mathbf{X}^{n+1}
  =
  \mathbf{X}^{n}
  +
  \mathbf{u}\,\delta t
  +
  \boldsymbol{\Sigma}\mathbf{W}_{\delta t}.
\end{equation}
It just needs the full factor \(\boldsymbol{\Sigma}\).
Footnote 4 sets \(\delta t=\Delta t\), which is exactly where the finite-step
cross covariance survives.

\section*{The issue}

One MC step chooses exactly one face, or stays put. So for \(i\neq j\),
\begin{equation}
  \mean{\Delta X_i\Delta X_j}=0.
\end{equation}
But that is the \emph{raw} cross moment, not the covariance:
\begin{align}
  \cov(\Delta X_i,\Delta X_j)
  &=
  \mean{\Delta X_i\Delta X_j}
  -\mean{\Delta X_i}\mean{\Delta X_j}
  \\
  &=
  -\mean{\Delta X_i}\mean{\Delta X_j}.
\end{align}
Therefore, whenever the mean transport has two nonzero components,
\begin{equation}
  \boxed{
  \cov(\Delta X_i,\Delta X_j)\neq0
  \qquad (i\neq j).
  }
\end{equation}

The 3D discussion calls
\(\mean{\Delta X_i\Delta X_j}/(2\Delta t)\) the covariance and then sets the
off-diagonal entries to zero because a single MC jump cannot move in two
directions. The second statement is right for the raw moment. The missing term
is the product of the means.

\section*{Full finite-step result}

For direction \(i\), let
\begin{align}
  p_i^+ &= P(\ddX=+h_i\mathbf{e}_i), &
  p_i^- &= P(\ddX=-h_i\mathbf{e}_i),\\
  C_{+,i} &= p_i^+ + p_i^-, &
  C_{-,i} &= p_i^+ - p_i^- .
\end{align}
The remaining probability is the no-jump outcome:
\begin{equation}
  P(\ddX=0)=1-\sum_i C_{+,i}.
\end{equation}

First moment:
\begin{equation}
  m_i\equiv\mean{\Delta X_i}=h_i C_{-,i}.
\end{equation}

Raw second moment:
\begin{equation}
  R_{ij}\equiv\mean{\Delta X_i\Delta X_j}
  =h_i^2 C_{+,i}\,\delta_{ij}.
\end{equation}

Central second moment:
\begin{align}
  Q_{ij}
  &\equiv
  \mean{(\Delta X_i-m_i)(\Delta X_j-m_j)}
  \\
  &=R_{ij}-m_i m_j.
\end{align}
So
\begin{equation}
  \boxed{
  Q_{ij}
  =
  h_i^2 C_{+,i}\,\delta_{ij}
  -
  h_i h_j C_{-,i}C_{-,j}.
  }
\end{equation}
Equivalently,
\begin{equation}
  \boxed{
  \mathbf{Q}
  =
  \diag\!\left(h_i^2 C_{+,i}\right)
  -
  \mathbf{m}\mathbf{m}^{\mathsf T}.
  }
\end{equation}

Diagonal:
\begin{equation}
  Q_{ii}=h_i^2\left(C_{+,i}-C_{-,i}^2\right),
\end{equation}
which is exactly the current 1D formula.

Off diagonal:
\begin{equation}
  \boxed{
  Q_{ij}=-h_i h_j C_{-,i}C_{-,j},
  \qquad i\neq j.
  }
\end{equation}

\section*{Why the infinitesimal argument looks diagonal}

For finite velocities,
\begin{equation}
  m_i=u_i\Delta t=O(\Delta t).
\end{equation}
Therefore,
\begin{equation}
  m_i m_j=O(\Delta t^2),
\end{equation}
and
\begin{equation}
  \frac{m_i m_j}{\Delta t}=O(\Delta t)\rightarrow0.
\end{equation}
So the off-diagonal Kramers--Moyal \emph{rate} vanishes as
\(\Delta t\rightarrow0\). No problem there.

The finite-step algorithm is different: it takes one increment with
\(\delta t=\Delta t\), and the diagonal coefficient already keeps the
finite-step correction \(-C_{-,i}^2\). The cross terms
\(-C_{-,i}C_{-,j}\) are the matching finite-step corrections off the diagonal.

\section*{Simple 2D example}

Uniform positive flow, equal cell widths, and
\begin{equation}
  C_x=C_y=C.
\end{equation}
The MC outcomes are
\begin{equation}
  \ddX=
  \begin{cases}
    (h,0), & P=C,\\
    (0,h), & P=C,\\
    (0,0), & P=1-2C.
  \end{cases}
\end{equation}
Then
\begin{align}
  \mean{\Delta X}
  &=
  \mean{\Delta Y}
  =hC,\\
  \operatorname{Var}(\Delta X)
  &=
  \operatorname{Var}(\Delta Y)
  =h^2C(1-C),\\
  \cov(\Delta X,\Delta Y)
  &=-h^2C^2,\\
  \operatorname{Corr}(\Delta X,\Delta Y)
  &=-\frac{C}{1-C}.
\end{align}

At \(C=0.1\),
\begin{equation}
  \operatorname{Corr}_{\rm MC}=-0.111111.
\end{equation}

AthenaK, \(131{,}072\) It\^o-2 particles:
\begin{center}
\begin{tabular}{@{}lrr@{}}
\toprule
 & measured & finite-step MC \\
\midrule
\(\operatorname{Var}(\Delta X)\)
  & \(8.7997\times10^{-5}\) & \(8.7891\times10^{-5}\)\\
\(\operatorname{Var}(\Delta Y)\)
  & \(8.7941\times10^{-5}\) & \(8.7891\times10^{-5}\)\\
\(\cov(\Delta X,\Delta Y)\)
  & \(-2.45\times10^{-8}\) & \(-9.77\times10^{-6}\)\\
\(\operatorname{Corr}(\Delta X,\Delta Y)\)
  & \(-2.78\times10^{-4}\) & \(-0.111111\)\\
\bottomrule
\end{tabular}
\end{center}

Exactly what independent coordinate kicks should do: right marginal
variances, zero cross covariance.

\section*{Proposed fix}

If the claim is only ``same mean and variance in each coordinate,'' the current
method is fine. I would just say that explicitly.

If the claim is ``same first two finite-step displacement moments,'' use the
full \(\mathbf{Q}\).

Factor
\begin{equation}
  \mathbf{Q}=\mathbf{L}\mathbf{L}^{\mathsf T},
\end{equation}
draw
\begin{equation}
  \mean{\boldsymbol{\xi}}=0,
  \qquad
  \mean{\boldsymbol{\xi}\boldsymbol{\xi}^{\mathsf T}}=\mathbf{I},
\end{equation}
and update
\begin{equation}
  \boxed{
  \ddX
  =
  \mathbf{m}
  +
  \mathbf{L}\boldsymbol{\xi}.
  }
\end{equation}
With the current notation,
\begin{equation}
  \mathbf{u}=\frac{\mathbf{m}}{\Delta t},
  \qquad
  \boldsymbol{\kappa}=\frac{\mathbf{Q}}{2\Delta t},
\end{equation}
so the same update is
\begin{equation}
  \ddX
  =
  \mathbf{u}\Delta t
  +
  \sqrt{2\Delta t}\,
  \boldsymbol{\kappa}^{1/2}\boldsymbol{\xi}.
\end{equation}

Bounded, independent, zero-mean, unit-variance components of
\(\boldsymbol{\xi}\) still work. The factor \(\mathbf{L}\) supplies the
correlation. This matches the MC mean and the complete finite-step covariance;
higher moments remain a separate question.

\medskip
\noindent\fbox{\parbox{0.94\linewidth}{
Bottom line: independent coordinate kicks get each coordinate's mean and
variance right, but not the full finite-step covariance. To match that too,
include \(-h_i h_j C_{-,i}C_{-,j}\) and draw the kick through
\(\mathbf{Q}=\mathbf{L}\mathbf{L}^{\mathsf T}\).

\textit{Am I reading that right?}\hfill Drummond
}}

\end{document}
